% =============================================================================
%  I. INTRODUCTION: THE CLINICAL PROBLEM          
% =============================================================================
\begin{frame}{I. Introduction: Accuracy Survives, Usefulness Does Not}
\begin{itemize}\itemsep2pt\small
  \item Sepsis kills about \textbf{eleven million people a year}, nearly one
        death in five worldwide~\cite{rudd2020global}; in the ICU every hour of
        delayed treatment raises the risk of death~\cite{seymour2017time}.
  \item Wang et al.~\cite{wang2025srpm} reviewed \textbf{91} real-time sepsis
        prediction models:
\end{itemize}

\begin{center}\tablestyle\footnotesize
\begin{tabular}{lcc}
\toprule
\textbf{Metric} & \textbf{Internal} & \textbf{External} \\
\midrule
Median AUROC         & 0.811            & 0.783 \\
Median Utility Score & \textbf{$+$0.381} & \textbf{$-$0.164} \\
\bottomrule
\end{tabular}
\end{center}

\begin{itemize}\itemsep2pt\small
  \item A negative Utility Score means the alerts cause more harm than benefit;
        AUROC did not detect the collapse, because it measures ranking, not
        whether alerts are timely or useful.
  \item The same gap appears in deployment: Epic Sepsis Model v2 across
        227,091 encounters, AUROC 0.82--0.92 with PPV
        0.13--0.26~\cite{wong2026epic}.
\end{itemize}

\vfill
\begin{alertblock}{}
  \centering ``Which model is best?'' is the wrong question to lead with.
\end{alertblock}
\end{frame}
